\documentclass[12pt,oneside]{book}

% ============================================================
% METADATOS
% ============================================================
\newcommand{\universidad}{Universidad de Buenos Aires (UBA)}
\newcommand{\facultad}{Facultad de Derecho}
\newcommand{\materia}{Derecho Civil}
\newcommand{\titulo}{Atributos de la Persona Humana: El Nombre y sus Reglas. La Filiación y el Apellido}
\newcommand{\autor}{Isaac Edgar Camacho Ocampo}
\newcommand{\anio}{2026}
\newcommand{\reflexionmotivacional}{Comprender los fundamentos del derecho civil es el primer paso para transformar el estudio en criterio jurídico.}

% ============================================================
% PAQUETES
% ============================================================
\usepackage[utf8]{inputenc}
\usepackage[T1]{fontenc}
\usepackage[spanish,es-nodecimaldot]{babel}

\usepackage[
    top=2.5cm,
    bottom=2.5cm,
    left=3cm,
    right=2.5cm,
    headheight=15pt
]{geometry}

\usepackage{amsmath}
\usepackage{amssymb}
\usepackage{mathtools}

\usepackage{graphicx}
\usepackage{float}
\usepackage{caption}
\usepackage{subcaption}

\usepackage{booktabs}
\usepackage{array}
\usepackage{longtable}
\usepackage{tabularx}

\usepackage{enumitem}

\usepackage{xcolor}
\usepackage[most]{tcolorbox}

\usepackage{fancyhdr}
\usepackage{ragged2e}

% ============================================================
% RUTAS DE IMÁGENES
% ============================================================
\graphicspath{
    {./img/}
    {./graficos/}
    {./figuras/}
}

% ============================================================
% CONFIGURACIÓN
% ============================================================
\setcounter{tocdepth}{2}

\setlength{\parindent}{0pt}
\setlength{\parskip}{0.5em}

\setlist[itemize]{
    topsep=0.3em,
    itemsep=0.2em
}
\setlist[enumerate]{
    topsep=0.3em,
    itemsep=0.2em
}

\clubpenalty=10000
\widowpenalty=10000

\pagestyle{fancy}
\fancyhf{}
\fancyhead[L]{\nouppercase{\leftmark}}
\fancyhead[R]{\materia}
\fancyfoot[C]{\thepage}
\renewcommand{\headrulewidth}{0.4pt}
\renewcommand{\footrulewidth}{0pt}

\fancypagestyle{plain}{
    \fancyhf{}
    \fancyfoot[C]{\thepage}
    \renewcommand{\headrulewidth}{0pt}
}

% ============================================================
% COMANDOS MATEMÁTICOS
% ============================================================
\newcommand{\abs}[1]{\left|#1\right|}
\newcommand{\norm}[1]{\left\lVert#1\right\rVert}
\newcommand{\vect}[1]{\mathbf{#1}}
\newcommand{\deriv}[2]{\frac{d#1}{d#2}}
\newcommand{\pderiv}[2]{\frac{\partial#1}{\partial#2}}

% ============================================================
% ENTORNOS / CAJAS ACADÉMICAS
% ============================================================
\newtcolorbox{definition}[1]{
    enhanced,
    breakable,
    title={Definición: #1},
    fonttitle=\bfseries,
    colback=blue!3,
    colframe=blue!50!black,
    boxrule=0.8pt,
    arc=2mm
}

\newtcolorbox{important}{
    enhanced,
    breakable,
    title={Importante},
    fonttitle=\bfseries,
    colback=yellow!8,
    colframe=orange!70!black,
    boxrule=0.8pt,
    arc=2mm
}

\newtcolorbox{example}[1]{
    enhanced,
    breakable,
    title={Ejemplo: #1},
    fonttitle=\bfseries,
    colback=green!4,
    colframe=green!40!black,
    boxrule=0.8pt,
    arc=2mm
}

\newtcolorbox{formula}[1]{
    enhanced,
    breakable,
    title={Regla legal: #1},
    fonttitle=\bfseries,
    colback=purple!4,
    colframe=purple!50!black,
    boxrule=0.8pt,
    arc=2mm
}

\newtcolorbox{exercise}[1]{
    enhanced,
    breakable,
    title={Ejercicio: #1},
    fonttitle=\bfseries,
    colback=gray!5,
    colframe=gray!60!black,
    boxrule=0.8pt,
    arc=2mm
}

\newtcolorbox{note}{
    enhanced,
    breakable,
    title={Nota},
    fonttitle=\bfseries,
    colback=white,
    colframe=gray!50,
    boxrule=0.6pt,
    arc=2mm
}

% ============================================================
% HIPERVÍNCULOS Y METADATOS PDF (al final, ya definidos los datos)
% ============================================================
\usepackage[
    colorlinks=true,
    linkcolor=blue!50!black,
    citecolor=green!40!black,
    urlcolor=blue!60!black,
    pdftitle={\titulo},
    pdfauthor={\autor},
    pdfsubject={Apuntes universitarios},
    bookmarks=true
]{hyperref}

% ============================================================
% DOCUMENTO
% ============================================================
\begin{document}

% ---------------- PORTADA ----------------
\begin{titlepage}
\thispagestyle{empty}

\begin{center}

\vspace*{1cm}

{\large \textsc{\universidad}}

\vspace{0.2cm}

{\large \textsc{\facultad}}

\vspace{3cm}

{\Huge \textbf{\materia}}

\vspace{1cm}

{\LARGE \titulo}

\vspace{0.8cm}

\textit{\reflexionmotivacional}

\vspace{1.2cm}

\rule{0.70\textwidth}{0.5pt}

\vspace{0.5cm}

{\large Apuntes de estudio}

\vspace{0.5cm}

\rule{0.70\textwidth}{0.5pt}

\vfill

{\large \autor}

\vspace{0.3cm}

{\large \anio}

\end{center}

\vfill

\begin{flushleft}
\textit{Este es un modesto aporte a los estudiantes de ``Derecho Civil'' no pretende sustituir el material oficial de la materia.}
\end{flushleft}

\end{titlepage}

% ---------------- ÍNDICE ----------------
\tableofcontents

% ============================================================
% CAPÍTULO 1 — BLOQUE I: ATRIBUTOS DE LA PERSONA HUMANA
% ============================================================
\chapter{Los atributos de la persona humana}

Tras haber estudiado la persona humana y el comienzo de su existencia, el derecho civil se ocupa de sus atributos, es decir, de los elementos que permiten que la personalidad jurídica se despliegue. Los atributos son los elementos que posibilitan que las personas, en tanto sujetos en relación, establezcan relaciones y situaciones jurídicas y puedan identificarse dentro de ellas.

\section{Enumeración de los atributos de la persona}

La distinción clásica de la doctrina reconoce cinco atributos de la persona humana: el nombre, el domicilio, el estado, la capacidad y el patrimonio.

\begin{table}[H]
\centering
\begin{tabularx}{\textwidth}{l X}
\toprule
\textbf{Atributo} & \textbf{Función} \\
\midrule
Nombre & Permite identificar a la persona. \\
Domicilio & Permite ubicar a la persona. \\
Estado & Permite establecer frente a qué relaciones jurídicas se encuentra la persona. \\
Capacidad & Permite determinar qué aptitud tiene la persona para entablar relaciones jurídicas. \\
Patrimonio & Permite conocer qué bienes tiene una persona. \\
\bottomrule
\end{tabularx}
\caption{Atributos de la persona humana y su función}
\end{table}

\begin{note}
Existen autores que discuten la clasificación clásica de los atributos de la persona, en particular en lo referido a la inclusión del patrimonio dentro de dicha categoría. No obstante, el desarrollo de la materia sigue en este punto la distinción clásica de nombre, domicilio, estado, capacidad y patrimonio.
\end{note}

\section{El nombre como atributo}

\begin{definition}{Nombre}
El nombre es la identificación exclusiva que corresponde a una persona.
\end{definition}

El nombre constituye un concepto intuitivo, en la medida en que toda persona se encuentra habitualmente familiarizada con la utilización de su propio nombre al completar formularios y realizar trámites cotidianos. Se trata, además, de una materia que se encuentra regulada de manera clara en el Código Civil y Comercial.

\textbf{Caracteres del nombre.} El nombre presenta los siguientes caracteres:

\begin{itemize}
\item \textbf{Necesario}
\item \textbf{Obligatorio}
\item \textbf{Único}
\item \textbf{Inalienable}
\item \textbf{Inembargable}
\item \textbf{Imprescriptible}
\item \textbf{Inmutable}
\item \textbf{Indivisible}
\end{itemize}

El nombre está compuesto por dos elementos:

\begin{itemize}
\item El \textbf{nombre individual}, también llamado nombre de pila o nombre familiar, que el Código denomina \textbf{prenombre}, y que es la denominación que permite distinguir a la persona dentro de su familia.
\item El \textbf{apellido}, que es la denominación común a los integrantes de una misma familia.
\end{itemize}

Ambos elementos configuran, en conjunto, la identificación de la persona: si bien coloquialmente se suele hablar de ``nombre y apellido'' por razones de practicidad, jurídica y técnicamente ambos integran el concepto único de \textbf{nombre}.

% ============================================================
% CAPÍTULO 2 — BLOQUE II: EL NOMBRE (ARTS. 62 A 72 CCyC)
% ============================================================
\chapter{El nombre}

\section{Régimen legal del nombre: evolución legislativa}

La regulación del nombre se ubica en el Código Civil y Comercial dentro del Libro Primero (Parte General), Título 1 (La persona humana), Capítulo 4 (El nombre), y comprende los artículos 62 a 72.

\begin{formula}{Ubicación normativa}
Código Civil y Comercial de la Nación. Libro Primero, Parte General, Título 1 (La persona humana), Capítulo 4 (El nombre), artículos 62 a 72.
\end{formula}

La regulación actual deroga la anterior Ley 18.248, conocida como la ``ley del nombre''.

\begin{table}[H]
\centering
\begin{tabularx}{\textwidth}{l X}
\toprule
\textbf{Etapa} & \textbf{Régimen} \\
\midrule
Código de Vélez (Ley 340) & Sin regulación específica del nombre; se regía por la costumbre. \\
Ley 18.248 (año 1969) & Primera ``ley del nombre'', complementaria al Código de Vélez. \\
Código Civil y Comercial vigente & Incorpora la regulación del nombre directamente en su texto (arts. 62 a 72), derogando la Ley 18.248. \\
\bottomrule
\end{tabularx}
\caption{Evolución legislativa de la regulación del nombre}
\end{table}

\section{El nombre como derecho y como deber (art. 62)}

\begin{formula}{Artículo 62, Código Civil y Comercial}
La persona humana tiene el derecho y el deber de usar el prenombre y el apellido que le corresponde.
\end{formula}

La doctrina mayoritaria sostiene que el uso del nombre y del apellido es, simultáneamente, un derecho y un deber. Es un \textbf{derecho subjetivo}, porque la persona tiene la prerrogativa de exigir que se le reconozca con un determinado nombre. Es también un \textbf{deber}, porque existe un interés público en identificar a las personas y en poder establecer su identidad. Como correlato de este deber, la persona tiene la obligación de identificarse con su propio nombre, sin poder usar un nombre falso ni mentir respecto de él.

\section{El prenombre: reglas y límites (art. 63)}

\subsection{Adquisición del prenombre y orden de elección}

\begin{formula}{Artículo 63, inciso a) --- Adquisición del prenombre}
El prenombre se adquiere por la inscripción en el acta de nacimiento, a pedido de quien ostente el derecho a elegirlo.
\end{formula}

El orden de prelación de quienes tienen el derecho de elegir el prenombre es el siguiente:

\begin{enumerate}
\item En primer lugar, corresponde a los \textbf{padres}, quienes pueden ejercer ese derecho en conjunto o bien uno solo de ellos, sobreentendiéndose en este último caso que cuenta con la autorización del otro (cuestión regulada, en lo referido a la responsabilidad parental, en el artículo 641 del Código).
\item Los padres pueden \textbf{delegar} la elección en un tercero, otorgándole la autorización correspondiente para que actúe como delegado.
\item Si uno de los padres no está presente, corresponde elegir al otro (por ejemplo, en el caso de una persona soltera).
\item Si no hay padre conocido (por ejemplo, en el caso de un niño abandonado), la elección la realizan los guardadores designados judicialmente, el Ministerio Público o el funcionario del Registro Civil y de Capacidad de las Personas.
\end{enumerate}

\subsection{Reglas y límites para la elección del prenombre}

La elección del prenombre parte de un principio de \textbf{libertad de elección}, que incluso permite inscribir a las personas con nombres aborígenes o derivados de voces aborígenes, autóctonas o latinoamericanas. A partir de este principio, el Código establece límites precisos:

\begin{important}
No pueden inscribirse más de tres prenombres.

No pueden inscribirse apellidos como prenombres.

No pueden inscribirse primeros prenombres idénticos a los primeros prenombres de hermanos vivos.

No pueden inscribirse prenombres extravagantes.
\end{important}

\begin{example}{Prohibición de usar apellidos como prenombres}
No es posible inscribir a un hijo con el prenombre ``Messi'' o ``Maradona'', dado que se trata de apellidos, y la ley prohíbe inscribir apellidos como si fueran prenombres.
\end{example}

La prohibición de repetir el primer prenombre entre hermanos vivos no se aplica a los casos de nombres compuestos, tales como:

\begin{example}{Nombres compuestos que exceptúan la regla de identidad entre hermanos}
María de las Mercedes, María del Rosario, María del Pilar.

Juan Bautista, Juan Segundo, Juan Manuel, Juan Ignacio.
\end{example}

Respecto de los nombres extravagantes, cabe señalar que la ley anterior a 1948 era más amplia, en cuanto disponía que no podían inscribirse nombres extravagantes o ridículos, que expresaran tendencias políticas o ideológicas, o que fueran equívocos respecto del sexo de quien los llevara. Todo ese conjunto de prohibiciones se sintetiza actualmente en la sola expresión ``extravagantes''.

\begin{example}{Alcance del concepto de nombre extravagante}
La calificación de un nombre como ``extravagante'' deja un margen abierto a discusión. Un nombre como ``Revolución'' probablemente sería considerado extravagante. En un sentido similar, en Estados Unidos existió en un momento la práctica de inscribir a los hijos con nombres de marcas famosas, como ``Amazon'' o ``Google'', bajo la idea de que el hijo contaría en el futuro con una marca que lo respaldara; este tipo de nombres también encuadraría dentro de la categoría de ``extravagante''.
\end{example}

\begin{note}
Quien define en la práctica si un nombre es o no extravagante es el Registro del Estado Civil. Si la persona interesada tiene un conflicto con esa decisión, puede recurrir a la vía judicial.
\end{note}

\section{El apellido según el tipo de filiación}

La regla general sobre el apellido surge del artículo 63 del Código, correspondiendo distinguir distintos supuestos según el tipo de filiación de la persona.

\subsection{Filiación matrimonial (art. 64)}

En este punto, el Código Civil y Comercial innovó respecto de la Ley 18.248, que establecía necesariamente que el hijo llevaba el apellido paterno.

\begin{formula}{Artículo 64 --- Apellido de los hijos matrimoniales}
El hijo matrimonial lleva el primer apellido de alguno de los cónyuges; en caso de no haber acuerdo sobre qué apellido llevará, se determina por sorteo realizado en el Registro del Estado Civil y Capacidad de las Personas.

A pedido de los padres, o del interesado con edad y madurez suficiente, se puede agregar el apellido del otro cónyuge.

Todos los hijos comunes deben llevar el mismo apellido y, en el mismo orden, que el primero de los hijos.
\end{formula}

El Código Civil y Comercial permite utilizar el apellido de cualquiera de los dos cónyuges. Si los padres no se ponen de acuerdo sobre cuál de los dos apellidos llevará el hijo, la cuestión se resuelve por sorteo en el propio Registro Civil, a pedido de los padres o del interesado con edad y madurez suficiente, quien también puede solicitar que se agregue el apellido del otro progenitor.

\begin{important}
El orden de apellidos que se determine para el primer hijo --por acuerdo de los padres o por sorteo-- queda fijado para todos los hijos comunes posteriores de ese matrimonio, quienes deben mantener el mismo orden de apellidos que el primero.
\end{important}

En cuanto a la posibilidad de agregar el apellido del otro progenitor a pedido del propio interesado, el Código utiliza el estándar de \textbf{edad y madurez suficiente}. Es el Registro Civil quien evalúa en concreto si la persona reúne dicha condición. Como criterio orientativo, se señala que la persona debería contar, como mínimo, con trece años, edad en la que el Código presume el discernimiento para los actos lícitos, sin perjuicio de que esta cuestión deba analizarse en cada caso concreto.

\subsection{Filiación extramatrimonial}

Cuando la filiación es extramatrimonial pueden presentarse tres supuestos distintos:

\begin{enumerate}
\item \textbf{Un solo progenitor reconoce al hijo} --por ejemplo, solo la madre, que es quien da a luz, mientras que el padre no se hace cargo--. En este caso, el hijo lleva el apellido de quien lo reconoce.
\item \textbf{El segundo progenitor reconoce al hijo con posterioridad al primero.} En este supuesto, los progenitores deben ponerse de acuerdo sobre el apellido; si no lo logran, la cuestión la decide el juez, teniendo en cuenta el interés superior del niño.
\item \textbf{Ambos progenitores reconocen al hijo al mismo tiempo.} En este caso se aplican las mismas reglas que en la filiación matrimonial: deben ponerse de acuerdo sobre el orden de apellidos y, si no lo logran, la cuestión se resuelve por sorteo.
\end{enumerate}

\subsection{Persona sin filiación determinada (art. 66)}

Puede ocurrir que una persona no tenga una filiación determinada. Esta circunstancia suele presentarse en zonas marginales, en familias sin contacto con instituciones estatales o de protección, en las que la persona nunca fue inscripta en el Registro Civil, aunque la comunidad ya la conozca por un determinado apellido. Existe una regla por la cual no se puede excluir a un niño de la educación escolar por no contar con DNI, lo que hace que, en la práctica, la persona pueda educarse sin haber sido formalmente inscripta.

En estos casos, la circunstancia debe ser anotada por el oficial del Registro Civil, utilizando el apellido que la persona ya esté usando, o bien asignándole un apellido común.

\begin{formula}{Artículo 66, Código Civil y Comercial}
La persona, con edad y grado de madurez suficiente, que carezca de apellido inscripto, puede pedir la inscripción del que está usando.
\end{formula}

\section{Apellido de los cónyuges, viudez y adopción}

\subsection{Apellido de los cónyuges durante el matrimonio}

El Código consagra la libertad de que, al casarse, cualquiera de los cónyuges pueda optar por usar el apellido del otro. Se trata de un derecho indistinto, que puede ejercer cualquiera de los dos, a diferencia del régimen anterior, en el cual era la mujer quien tomaba el apellido del marido. El Código actual está redactado en términos neutros.

\begin{formula}{Régimen del apellido conyugal}
Cualquiera de los cónyuges puede optar por usar el apellido del otro, con la partícula ``de'' o sin ella, como apellido adicional al propio.
\end{formula}

\subsection{Apellido conyugal y divorcio}

Si los cónyuges se divorcian, quien adoptó el apellido del otro deja de usarlo. La misma regla se aplica en los casos de nulidad del matrimonio, salvo que, por motivos razonables, un juez autorice al cónyuge a conservar el apellido.

\begin{example}{Motivo razonable para conservar el apellido tras el divorcio: uso comercial}
Un cónyuge que utilizaba el apellido adquirido con el matrimonio también para identificar a su empresa --era conocido en el mercado con el nombre de la empresa según ese apellido-- continuó usándolo tras el divorcio por razones económicas, dado que era el nombre con el que su empresa era reconocida.
\end{example}

\subsection{Apellido del cónyuge viudo o viuda}

\begin{formula}{Régimen del apellido en caso de viudez}
El cónyuge viudo o viuda puede seguir usando el apellido del otro cónyuge, mientras no contraiga nuevas nupcias ni constituya una nueva unión convivencial. El uso del apellido no es obligatorio.
\end{formula}

\subsection{Apellido del hijo adoptado}

Según el tipo de adopción, existen variantes en el régimen del apellido del hijo adoptado. La adopción puede ser plena, simple o de integración, y cada una sigue reglas distintas.

\paragraph{Adopción plena.} Si la adopción es plena, el adoptado lleva el apellido del adoptante. Si la adopción es conjunta (por dos adoptantes), se aplican las mismas reglas que en la filiación matrimonial respecto del orden de los apellidos.

\begin{important}
Excepcionalmente, fundado en el derecho a la identidad, y a pedido de la parte interesada, se puede solicitar agregar o anteponer el apellido de origen del adoptado. La regla general en la adopción plena es tomar el apellido del adoptante, pero esta excepción busca preservar, en la medida de lo posible, el derecho a la identidad del niño respecto de su origen, en su dimensión estática.
\end{important}

\paragraph{Adopción simple.} La adopción simple no corta el vínculo con la familia de origen; no es tan radical como la adopción plena. Aquí la regla es la inversa respecto de la plena:

\begin{formula}{Régimen del apellido en la adopción simple}
El adoptado puede solicitar mantener su apellido de origen, adicionándole el del adoptante. Si no manifiesta una voluntad expresa al respecto, se aplica la misma regla que en la adopción plena en cuanto a tomar el apellido del adoptante (o de los adoptantes). A diferencia de la adopción plena, en la adopción simple la consulta previa al adoptado sobre mantener su apellido de origen es la regla, mientras que en la plena esa posibilidad es excepcional.
\end{formula}

\paragraph{Adopción de integración.} Se produce cuando se adopta al hijo del cónyuge o conviviente. En este caso se aplican las mismas reglas que correspondan según el tipo de adopción de que se trate (plena o simple).

\begin{note}
El artículo 64 solo permite agregar apellidos, no suprimirlos. Cabe distinguir, además, entre \textbf{prenombre} (nombre individual o nombre de pila) y \textbf{nombre} en sentido amplio y técnico, que está compuesto por dos partes: el prenombre y el apellido. Coloquialmente se suele decir ``nombre y apellido'' por razones de practicidad, pero jurídica y técnicamente ambos integran el concepto único de nombre.

Al prenombre también se lo conoce como ``nombre de pila'', denominación que obedece a la pila bautismal, tradición de raíz cristiana extendida posteriormente por todo el mundo. El Código actual, no obstante esa tradición histórica, utiliza la palabra técnica ``prenombre''.
\end{note}

\section{Cambio de nombre y acciones de protección}

\subsection{Principio de inmutabilidad}

\begin{important}
El nombre es inmutable. El cambio de nombre y apellido solo procede si existen justos motivos.
\end{important}

El cambio de nombre puede ocurrir por dos vías: una vía administrativa y una vía judicial.

\subsection{Vía judicial: justos motivos del art. 69}

En la vía judicial es necesario alegar y acreditar un justo motivo. El artículo 69 del Código enumera algunos justos motivos a título orientativo:

\begin{formula}{Artículo 69 --- Justos motivos para el cambio de nombre por vía judicial (enumeración no taxativa)}
El seudónimo, cuando hubiere adquirido notoriedad.

La raigambre cultural, étnica o religiosa.

La afectación de la personalidad de la persona interesada, cualquiera sea su causa, siempre que se encuentre acreditada.
\end{formula}

Respecto del \textbf{seudónimo notorio}, esta situación se presenta habitualmente entre artistas que, con el tiempo, adquieren notoriedad a través de un seudónimo. La \textbf{raigambre cultural, étnica o religiosa} se refiere al caso de una persona que descubre una raíz étnica o religiosa propia y desea reflejar ese origen en su nombre. La \textbf{afectación de la personalidad} se vincula con nombres que han provocado situaciones de bullying, burlas o vergüenza social, motivando el pedido de cambio.

\begin{note}
Estos tres motivos no son los únicos posibles: se trata de una lista enunciativa. Puede haber otros justos motivos, siempre que tengan la fuerza necesaria como para convencer a un juez de que corresponde cambiar el nombre, en el marco del proceso judicial específico previsto a tal efecto.
\end{note}

\subsection{Vía administrativa}

La vía administrativa, también prevista en el artículo 69, se aplica en dos supuestos en los cuales la ley autoriza directamente el cambio de nombre, sin necesidad de recurrir a un juez.

\begin{formula}{Supuestos de cambio de nombre por vía administrativa (art. 69)}
Por razón de identidad de género, conforme a la Ley 26.743.

Por haber sido víctima de desaparición forzada, sustitución de identidad, alteración o supresión del estado civil, o de un delito que afecte la personalidad.
\end{formula}

En estos dos supuestos se recurre directamente al Registro Civil, se acreditan los recaudos correspondientes y se obtiene el cambio de nombre sin necesidad de un proceso judicial.

\subsection{Procedimiento y acciones de protección del nombre}

El procedimiento para el cambio de nombre está regulado en el artículo 70, mientras que el artículo 71 regula las acciones vinculadas a la protección del nombre.

\begin{formula}{Artículo 71 --- Acciones de protección del nombre}
Acción de reclamación: corresponde cuando a una persona se le desconoce el uso de su propio nombre.

Acción de impugnación o usurpación del nombre: corresponde cuando otra persona usa el nombre de quien reclama, causándole un perjuicio.

Acción de defensa del nombre (supresión del nombre): corresponde cuando el uso del nombre por un tercero, por ejemplo como nombre de fantasía en una serie o telenovela, causa un perjuicio a quien lleva ese nombre.
\end{formula}

\begin{example}{Acción de reclamación}
Si una universidad anota incorrectamente el nombre de una persona, esta puede reclamar y solicitar que se lo anote correctamente.
\end{example}

\begin{example}{Acción de impugnación por usurpación}
Si una persona abre una cuenta en una red social con el nombre de otra y comienza a actuar en su nombre, la persona afectada puede iniciar una acción judicial para que se deje de usar su nombre.
\end{example}

\begin{example}{Acción de defensa o supresión del nombre}
Si el nombre de una persona aparece como nombre de fantasía en una serie o telenovela y ello le causa un perjuicio, puede pedir judicialmente que se deje de usar su nombre en ese contexto.
\end{example}

\subsection{El seudónimo}

\begin{definition}{Seudónimo}
Nombre ficticio empleado por una persona, generalmente en el ámbito artístico, en el ámbito de despliegue de su actividad. El uso del seudónimo no es ilícito; por el contrario, es lícito.
\end{definition}

\begin{formula}{Régimen del seudónimo}
Si el seudónimo goza de notoriedad, tiene la misma protección que el nombre, y le corresponden las mismas acciones previstas en el artículo 71 para la protección del nombre.
\end{formula}

% ============================================================
% CAPÍTULO 3 — BLOQUE III: LA FILIACIÓN
% ============================================================
\chapter{La filiación}

\section{Concepto de filiación}

La filiación es un concepto jurídico: es el estatuto de hijo, es decir, tener el estado de hijo. Una persona es hija y, por lo tanto, tiene una filiación, entendida como un vínculo entre un progenitor y un hijo.

\begin{definition}{Filiación}
La filiación es un vínculo jurídico entre progenitores e hijos, que tiene su origen ya sea en la naturaleza, ya sea en un reconocimiento jurídico por el cual una persona reconoce a su hijo extramatrimonial, ya sea en la fecundación in vitro (que tiene su propia regulación dentro del derecho de familia), o ya sea en la adopción.
\end{definition}

\section{Filiación matrimonial y su fundamento}

La filiación puede determinarse matrimonialmente: cuando dos personas se casan, los hijos que tengan se presumen del matrimonio. Si bien esta presunción puede no coincidir efectivamente con la realidad biológica en todos los casos (por ejemplo, ante una infidelidad), en principio la ley presume que la persona nacida es hija de ese matrimonio.

\begin{important}
La ley establece esta presunción por razones de orden social, en tanto el matrimonio es ordenador de la vida social: los cónyuges se comprometen a vivir juntos y a educar a los hijos que nazcan, y la sociedad, a través de la ley, reconoce ese compromiso otorgando a esos hijos el apellido fijado y presumido legalmente.
\end{important}

\section{Filiación por naturaleza y modelo binario legal}

Al analizar el caso de la filiación extramatrimonial cuando existen dos vínculos filiales --una madre y un padre-- corresponde limitarse, a los fines de simplificar la explicación, al caso de la filiación por naturaleza, es decir, sin intervención de técnicas de reproducción humana asistida.

\begin{formula}{Determinación de la maternidad y la paternidad}
Cuando una mujer da a luz biológicamente a un hijo, se anota la filiación materna: la maternidad queda determinada por el parto.

En cuanto a la paternidad, si la madre está casada, la ley presume la paternidad del cónyuge de la madre. Si la mujer no está casada, la persona tiene, en principio, un solo progenitor determinado --la madre--, salvo que otra persona reconozca al hijo.
\end{formula}

\begin{note}
En estos casos, la ley se basa en un modelo binario que sigue a la biología, en tanto se trata de filiación por naturaleza, distinta de otras formas de filiación que el propio Código también regula.
\end{note}

\section{Filiación extramatrimonial y acción de reclamación}

La filiación extramatrimonial puede ser fruto de un reconocimiento voluntario o de una acción judicial, cuando el padre no reconoce al hijo de manera espontánea.

\begin{definition}{Acción de filiación forzosa}
Acción judicial mediante la cual se obliga al padre que no reconoce voluntariamente a su hijo a reconocerlo y a pasar alimentos. En el momento en que se asigna el vínculo filial, el hijo pasa a tener determinados a ambos progenitores, que son quienes biológicamente dieron origen a ese nacimiento.
\end{definition}

\begin{example}{Reconocimiento como acto jurídico}
El reconocimiento voluntario de un hijo extramatrimonial por parte de un progenitor constituye, técnicamente, un acto jurídico de reconocimiento.
\end{example}

\section{Fuentes de la filiación}

Conforme la definición expuesta, la filiación puede tener su origen en la naturaleza, en el reconocimiento jurídico voluntario, en las técnicas de reproducción humana asistida, o en la adopción. Estas son las distintas formas que puede adoptar la filiación determinada, a diferencia de los casos de filiación no determinada, en los que se desconoce el vínculo filial y, conforme lo explicado respecto del artículo 66, se utiliza el apellido que ya esté empleando el hijo.

\section{El interés superior del niño}

\begin{formula}{Referencias normativas sobre el interés superior del niño}
Convención sobre los Derechos del Niño.

Ley 26.061 (ley especial sobre protección integral del interés superior del niño).
\end{formula}

Cuando el juez debe fallar teniendo en cuenta el interés superior del niño --por ejemplo, ante un desacuerdo sobre el apellido de un hijo extramatrimonial reconocido sucesivamente por ambos progenitores-- se remite centralmente a la Convención sobre los Derechos del Niño, sin perjuicio de la existencia de la Ley 26.061 como ley especial sobre la materia.

El interés superior del niño es un concepto complejo que presenta un componente objetivo --determinar qué es lo mejor para ese niño en concreto-- y que también tiene en cuenta la opinión del propio niño, de acuerdo a su edad y grado de madurez. En definitiva, se ponderan un conjunto amplio de elementos.

% ============================================================
% CUADRO CORNELL DE REPASO
% ============================================================
\chapter{Cuadro Cornell de repaso}

La siguiente tabla organiza, en formato Cornell, los conceptos centrales desarrollados a lo largo del presente material: a la izquierda, las preguntas o conceptos clave; a la derecha, su desarrollo o respuesta. Al final se incluye una síntesis integradora de todo el contenido.

\begin{longtable}{
    >{\raggedright\arraybackslash}p{0.32\textwidth}
    >{\raggedright\arraybackslash}p{0.58\textwidth}
}
\caption{Cuadro Cornell --- El nombre y la filiación} \\
\toprule
\textbf{Preguntas / palabras clave} & \textbf{Notas / respuestas} \\
\midrule
\endfirsthead

\multicolumn{2}{c}{\textbf{\tablename\ \thetable{} (continuación)}} \\
\toprule
\textbf{Preguntas / palabras clave} & \textbf{Notas / respuestas} \\
\midrule
\endhead

\midrule
\multicolumn{2}{r}{\textit{Continúa en la página siguiente}} \\
\endfoot

\bottomrule
\endlastfoot

\textbf{¿Cuáles son los atributos de la persona humana?} &
Nombre, domicilio, estado, capacidad y patrimonio: elementos que permiten que la personalidad jurídica se despliegue y que la persona establezca relaciones y situaciones jurídicas. \\

\textbf{¿Qué es el nombre?} &
La identificación exclusiva que corresponde a una persona. Es necesario, obligatorio, único, inalienable, inembargable, imprescriptible, inmutable e indivisible. Está compuesto por el prenombre y el apellido. \\

\textbf{¿Dónde está regulado el nombre?} &
En el Código Civil y Comercial, Libro Primero, Parte General, Título 1 (La persona humana), Capítulo 4, artículos 62 a 72. Deroga la Ley 18.248 (de 1969), que antes regulaba el nombre de forma complementaria al Código de Vélez. \\

\textbf{¿El nombre es un derecho o un deber?} &
Ambas cosas. Es un derecho subjetivo, porque la persona puede exigir que se la reconozca con su nombre; y es un deber, porque existe un interés público en identificar a las personas, lo que impide usar un nombre falso. \\

\textbf{¿Quién elige el prenombre de un recién nacido?} &
En primer lugar, los padres (en conjunto o uno solo con autorización del otro), quienes pueden delegar en un tercero; si falta un progenitor, decide el otro; y si no hay padre conocido, deciden los guardadores judiciales, el Ministerio Público o el funcionario del Registro Civil. \\

\textbf{¿Qué límites existen para elegir el prenombre?} &
No pueden inscribirse más de tres prenombres; no pueden inscribirse apellidos como prenombres; no pueden repetirse primeros prenombres entre hermanos vivos (salvo nombres compuestos); y no pueden inscribirse prenombres extravagantes. \\

\textbf{¿Qué apellido lleva el hijo matrimonial?} &
El primer apellido de cualquiera de los cónyuges, según acuerdo entre ellos; si no hay acuerdo, se decide por sorteo en el Registro Civil. A pedido de los padres o del interesado con edad y madurez suficiente, puede agregarse el apellido del otro cónyuge. Todos los hijos comunes llevan el mismo orden de apellidos que el primero. \\

\textbf{¿Qué apellido lleva el hijo extramatrimonial?} &
Si solo un progenitor lo reconoce, lleva su apellido. Si el segundo lo reconoce después, deben acordar el apellido y, si no acuerdan, decide el juez según el interés superior del niño. Si ambos lo reconocen simultáneamente, se aplican las reglas de la filiación matrimonial. \\

\textbf{¿Qué pasa si una persona no tiene filiación determinada?} &
El oficial del Registro Civil anota a la persona con el apellido que ya esté usando, o le asigna un apellido común. Además, según el artículo 66, la propia persona con edad y madurez suficiente puede pedir la inscripción del apellido que está usando. \\

\textbf{¿Puede un cónyuge usar el apellido del otro?} &
Sí. Cualquiera de los cónyuges puede optar por usar el apellido del otro (con ``de'' o sin ella) como apellido adicional. Es una opción neutra, sin distinción de género, a diferencia del régimen anterior. \\

\textbf{¿Qué ocurre con el apellido marital tras el divorcio?} &
Por regla general, se deja de usar el apellido del otro cónyuge, salvo que, por motivos razonables (por ejemplo, su uso comercial o empresarial consolidado), un juez autorice a conservarlo. \\

\textbf{¿Qué ocurre con el apellido marital en caso de viudez?} &
El cónyuge viudo o viuda puede seguir usando el apellido del otro cónyuge, de manera no obligatoria, mientras no contraiga nuevas nupcias ni constituya una nueva unión convivencial. \\

\textbf{¿Qué apellido lleva el hijo adoptado por adopción plena?} &
Lleva el apellido del adoptante (o de ambos adoptantes, según las reglas de la filiación matrimonial si son dos). Excepcionalmente, fundado en el derecho a la identidad, puede solicitarse agregar o anteponer el apellido de origen. \\

\textbf{¿Qué apellido lleva el hijo adoptado por adopción simple?} &
Puede solicitar mantener su apellido de origen, adicionándole el del adoptante. Si no manifiesta nada, se aplica la misma regla que en la adopción plena, pero aquí la consulta previa al adoptado es la regla, no la excepción. \\

\textbf{¿Qué reglas rigen la adopción de integración?} &
Se aplican las mismas reglas de apellido que correspondan según el tipo de adopción (plena o simple), ya que la adopción de integración es la adopción del hijo del cónyuge o conviviente. \\

\textbf{¿Qué es la filiación?} &
Un vínculo jurídico entre progenitores e hijos, que da el estatuto o estado de hijo. Su origen puede ser la naturaleza, el reconocimiento jurídico voluntario, las técnicas de reproducción humana asistida, o la adopción. \\

\textbf{¿Cómo se determina la maternidad y la paternidad por naturaleza?} &
La maternidad queda determinada por el parto. La paternidad se presume respecto del cónyuge de la madre si esta está casada; si no lo está, la madre es la única progenitora determinada, salvo que otra persona reconozca al hijo. \\

\textbf{¿Qué es la acción de filiación forzosa?} &
La acción judicial mediante la cual se obliga a un progenitor que no reconoce voluntariamente a su hijo a reconocerlo y a asumir sus obligaciones, como la de alimentos. \\

\textbf{¿Cuál es el principio general sobre el nombre?} &
El nombre es inmutable. Solo puede cambiarse si existen justos motivos, ya sea por vía judicial o por vía administrativa. \\

\textbf{¿Cuáles son los justos motivos previstos en el artículo 69 para la vía judicial?} &
El seudónimo notorio, la raigambre cultural, étnica o religiosa, y la afectación acreditada de la personalidad (por ejemplo, por bullying). La lista es enunciativa, no taxativa. \\

\textbf{¿En qué casos el cambio de nombre se tramita por vía administrativa?} &
Por razón de identidad de género (Ley 26.743), y por haber sido víctima de desaparición forzada, sustitución, alteración o supresión de identidad, o de un delito que afecte la personalidad. En estos casos se va directamente al Registro Civil. \\

\textbf{¿Qué acciones protegen el nombre (artículo 71)?} &
La acción de reclamación (si se desconoce el nombre de una persona), la acción de impugnación o usurpación (si otro usa el nombre sin derecho) y la acción de defensa o supresión del nombre (si su uso por un tercero causa perjuicio). \\

\textbf{¿Qué es el seudónimo y qué protección tiene?} &
Un nombre ficticio, generalmente de uso artístico, que es lícito. Si adquiere notoriedad, goza de la misma protección legal que el nombre, con las mismas acciones del artículo 71. \\

\midrule
\multicolumn{2}{p{0.90\textwidth}}{
\textbf{Resumen:} La clase desarrolló el nombre como uno de los cinco atributos de la persona humana, definiéndolo como la identificación exclusiva de cada persona, regulada en los artículos 62 a 72 del Código Civil y Comercial, que reemplazaron a la antigua Ley 18.248. El nombre se compone de dos elementos: el prenombre (de elección libre, pero limitada por reglas precisas: máximo tres nombres, prohibición de usar apellidos, de repetir nombres entre hermanos y de elegir nombres extravagantes) y el apellido, cuya asignación depende centralmente del tipo de filiación de la persona: matrimonial (elección entre cónyuges o sorteo), extramatrimonial (según el orden y simultaneidad del reconocimiento de los progenitores), por adopción (plena o simple, con distinto grado de protección del derecho a la identidad de origen) o sin filiación determinada.

A esto se suman reglas específicas sobre el apellido conyugal (opcional, reversible tras el divorcio salvo autorización judicial, y conservable en la viudez salvo nuevas nupcias). Finalmente, pese al principio de inmutabilidad del nombre, este puede cambiarse por justos motivos, ya sea por vía judicial (seudónimo notorio, raigambre cultural o étnica, afectación de la personalidad, entre otros motivos no taxativos) o por vía administrativa directa (identidad de género, víctimas de delitos graves contra la identidad), y el ordenamiento prevé acciones específicas (reclamación, impugnación y defensa del nombre) para proteger este atributo esencial de la persona, extendiendo dicha protección también al seudónimo notorio.

En definitiva, el nombre, lejos de ser un dato administrativo menor, es la expresión jurídica de la identidad de cada persona, profundamente entrelazada con el derecho de familia y con el derecho a la identidad reconocido constitucionalmente.
} \\

\end{longtable}

% ============================================================
% NORMATIVA CITADA
% ============================================================
\chapter*{Normativa citada}
\addcontentsline{toc}{chapter}{Normativa citada}

A continuación se listan las principales fuentes normativas citadas a lo largo del presente material:

\begin{itemize}
\item Código Civil y Comercial de la Nación, Libro Primero, Parte General, Título 1 (La persona humana), Capítulo 4 (El nombre), artículos 62 a 72.
\item Código Civil y Comercial de la Nación, artículo 641 (responsabilidad parental, referido al ejercicio conjunto o unilateral de la elección del prenombre).
\item Ley 18.248 (1969), ley del nombre, derogada por el actual Código Civil y Comercial.
\item Ley 26.743, de identidad de género.
\item Ley 26.061, de protección integral de los derechos de niñas, niños y adolescentes.
\item Convención sobre los Derechos del Niño.
\end{itemize}

\end{document}
